\input{preamble}

\title{\GlobalTitle}
\author{Киселев Николай, Б05-424}

\begin{document}
\maketitle

\begin{abstract}
    В данном проекте рассматривается задача \(\mathsf{MCDS}\) (minimum conntected dominating set --- минимальное связное доминирующее множество). Я разоберу связь данной задачи и \(\mathsf{MLST}\) (maximum leaf spanning tree --- остов с максимальным количеством листьев), приведу реализацию алгоритма её решения за \(O(1.9407^n)\), а также рассмотрю другие исследования, связанные с ней.
\end{abstract} 

\section{Введение}

На сегодняшний день, математическое сообщество считает, что \(\mathbf{P} \ne \mathbf{NP}\). Этот вопрос связан с очень многими областями в данной науке, например, с криптографией и кибербезопасностью в целом. Многие алгоритмы шифрования (например, RSA, DHE) используют неравенство двух классов для корректной работы, в частности, (предполагаемую) \(\mathbf{NP}\)-полноту задач \(\mathsf{FACTORING}\) --- разложения на множители и \(\mathbf{DLP}\) --- задача дискретного логарифмирования.

Данный проект детально изучит (предполагаемо) \(\mathbf{NP}\)-полную задачу (с доказательством принадлежности \(\mathbf{NPC}\)) \(\mathsf{MCDS}\) и некоторые ее взаимосвязи с другими похожими задачами (\(\mathsf{MLST}\)). Эта тема релевантна, как это было показано в \cite{wu}: данная проблема имеет большое значение для построения беспроводных ad-hoc сетей. Минимальное доминирующее множество в них обычно является множеством "шлюзов": для отправки пакета от хоста к хосту, он пересылается шлюзу (здесь используется свойство доминирующего множества), дальше происходит маршрутизация внутри доминирующего подграфа, после чего данные отправляются получателю.

Реализацию алгоритма я возьму из работы \cite{fomin}: в ней приведена оценка на \(O(1.9407^n)\). Стоить отметить также, что имеются и другие, более точные, например в \cite{fernau} было показано, что существует программа, решающая \(\mathsf{MLST}\) за \(O(1.8966^n)\). Это позволяет дать \(\mathsf{MCDS}\) оценку \(O\left((1.8966c)^n\right)\) для любого \(c > 1\) (будет рассмотрено далее в работе).

\printbibliography

\end{document} 
